\documentclass[12pt,a4paper]{amsart}

\usepackage{amsmath,amssymb,amsthm}
\usepackage{mathtools}
\usepackage[utf8]{inputenc}
\usepackage[ngerman]{babel}
\usepackage{hyperref}
\usepackage{geometry}
\geometry{margin=2.5cm}
\usepackage{booktabs}

% -------------------------------------------------------
\newtheorem{theorem}{Satz}[section]
\newtheorem{lemma}[theorem]{Lemma}
\newtheorem{korollar}[theorem]{Korollar}
\theoremstyle{definition}
\newtheorem{definition}[theorem]{Definition}
\newtheorem{bemerkung}[theorem]{Bemerkung}
\newtheorem{vermutung}[theorem]{Vermutung}
\newtheorem{beispiel}[theorem]{Beispiel}

% -------------------------------------------------------
\newcommand{\N}{\mathbb{N}}
\newcommand{\Z}{\mathbb{Z}}
\newcommand{\Q}{\mathbb{Q}}
\newcommand{\R}{\mathbb{R}}
\newcommand{\C}{\mathbb{C}}
\newcommand{\eq}[1]{e\!\left(#1\right)}

% -------------------------------------------------------
\begin{document}

\title{Die Goldbach'sche singuläre Reihe und die\\
       Hardy--Littlewood-Vermutung~$\mathbf{H}$\\
       über Primzahlkonstellationen}

\author{Michael Fuhrmann}
\address{Specialist Mathematics Project, Build 80}
\email{specialist-maths@localhost}
\date{2026}

\begin{abstract}
Die Hardy--Littlewood-Vermutungen (1923) liefern präzise quantitative
Vorhersagen für additive Probleme mit Primzahlen.  Für das binäre Goldbach-Problem
lautet die vorhergesagte gewichtete Darstellungsanzahl einer großen geraden
Zahl $n$ als $n = p + q$:
\[
  R_2(n) \;\sim\; \mathfrak{S}(n) \cdot n,
\]
wobei die \emph{Goldbach'sche singuläre Reihe}
\[
  \mathfrak{S}(n)
  \;=\; 2 C_2 \prod_{\substack{p \mid n \\ p > 2}} \frac{p-1}{p-2},
  \qquad
  C_2 \;=\; \prod_{p \geq 3} \frac{p(p-2)}{(p-1)^2} \approx 0{,}6602
\]
die lokale Arithmetik von $n$ kodiert.  Wir leiten diese Formel her,
beweisen $\mathfrak{S}(n) > 0$ für alle geraden $n \geq 4$, erklären
das Paritätsproblem und formulieren die Hardy--Littlewood-Vermutung~$\mathbf{H}$
über Primzahlkonstellationen ($k$-Tupel), die alle Primzahlzwillings-
und Goldbach-Fragen vereinheitlicht.
\end{abstract}

\maketitle
\tableofcontents

% ================================================================
\section{Einleitung}
% ================================================================

Das binäre Goldbach-Problem — jede gerade Zahl $n \geq 4$ ist Summe
zweier Primzahlen — ist numerisch bis $4 \times 10^{18}$ verifiziert
(Oliveira e Silva et al., 2014), aber unbewiesen.  Im Gegensatz zum
ternären Goldbach-Problem (bewiesen von Helfgott 2013, Paper~18) liegt
das binäre Problem außerhalb der Reichweite der Kreismethode allein,
wegen des Paritätsproblems.

Dennoch ist die \emph{erwartete} Darstellungsanzahl gut verstanden.
Definiere die gewichtete Darstellungsanzahl:
\[
  R_2(n) \;=\; \sum_{\substack{p_1 + p_2 = n}} \log p_1 \cdot \log p_2
  \;=\; \int_0^1 S(\alpha)^2\, \eq{-n\alpha}\, d\alpha,
\]
wobei $S(\alpha) = \sum_{m \leq n} \Lambda(m) \eq{m\alpha}$.

Die Kreismethode liefert:

\begin{vermutung}[Hardy--Littlewood, 1923]\label{verm:HL}
  Für jede gerade Zahl $n \geq 4$:
  \[
    R_2(n) \;\sim\; \mathfrak{S}(n) \cdot n,
  \]
  bzw.\ äquivalent: die Darstellungsanzahl $r_2(n) = \#\{p : p, n-p \in \mathbb{P}\}$
  erfüllt $r_2(n) \sim \mathfrak{S}(n) \cdot n / (\log n)^2$.
\end{vermutung}

% ================================================================
\section{Die Goldbach'sche singuläre Reihe}
\label{sec:singulaerreihe}
% ================================================================

\subsection{Herleitung durch die Kreismethode}

Aus der Hauptbogen-Approximation $S(\alpha) \approx \frac{\mu(q)}{\phi(q)} I(\beta)$
(Lemma~3.1 von Paper~18) folgt formal:

\begin{definition}[Goldbach'sche singuläre Reihe]\label{def:sing}
  \[
    \mathfrak{S}(n)
    \;=\; \sum_{q=1}^{\infty} \frac{\mu(q)^2\, c_q(n)}{\phi(q)^2},
  \]
  wobei $c_q(n) = \sum_{\gcd(a,q)=1} \eq{an/q}$ die Ramanujan-Summe ist.
\end{definition}

\subsection{Euler-Produktdarstellung}

\begin{theorem}[Euler-Produkt für $\mathfrak{S}(n)$]\label{thm:euler}
  Für gerades $n \geq 2$:
  \[
    \mathfrak{S}(n)
    \;=\; 2 C_2 \prod_{\substack{p \mid n \\ p \geq 3}} \frac{p-1}{p-2},
    \quad
    C_2 \;=\; \prod_{p \geq 3} \frac{p(p-2)}{(p-1)^2}.
  \]
\end{theorem}

\begin{proof}
  Die Ramanujan-Summe ist multiplikativ: $c_q(n) = \prod_{p|q} c_p(n)$ für
  quadratfreie $q$.  Da $\mu(q)^2 = 1$ nur für quadratfreie $q$, faktorisiert
  die Summe als Euler-Produkt:
  \[
    \mathfrak{S}(n) = \prod_p \left(1 + \frac{c_p(n)}{(p-1)^2}\right).
  \]
  \textit{Faktor bei $p=2$ (gerades $n$):}
  $c_2(n) = (-1)^n = 1$ (da $n$ gerade), also Faktor $= 1+1/1^2=2$.

  \textit{Faktor bei $p \geq 3$ mit $p \mid n$:}
  $c_p(n) = p-1$, Faktor $= 1+(p-1)/(p-1)^2 = 1+1/(p-1) = p/(p-1)$.

  \textit{Faktor bei $p \geq 3$ mit $p \nmid n$:}
  $c_p(n) = -1$, Faktor $= 1 - 1/(p-1)^2 = p(p-2)/(p-1)^2$.

  Daher:
  \begin{align*}
    \mathfrak{S}(n)
    &= 2 \cdot \prod_{\substack{p \geq 3 \\ p \mid n}} \frac{p}{p-1}
        \cdot \prod_{\substack{p \geq 3 \\ p \nmid n}} \frac{p(p-2)}{(p-1)^2} \\
    &= 2 \cdot \underbrace{\prod_{p \geq 3} \frac{p(p-2)}{(p-1)^2}}_{= C_2}
        \cdot \prod_{\substack{p \geq 3 \\ p \mid n}} \frac{p/(p-1)}{p(p-2)/(p-1)^2}
     = 2 C_2 \prod_{\substack{p \geq 3 \\ p \mid n}} \frac{p-1}{p-2}.  \qedhere
  \end{align*}
\end{proof}

% ================================================================
\section{Positivität von \texorpdfstring{$\mathfrak{S}(n)$}{S(n)}}
\label{sec:pos}
% ================================================================

\begin{theorem}\label{thm:pos}
  Für jede gerade Zahl $n \geq 4$: $\mathfrak{S}(n) > 0$.
\end{theorem}

\begin{proof}
  Alle Faktoren in Satz~\ref{thm:euler} sind positiv:
  \begin{itemize}
    \item $C_2 = \prod_{p\geq 3} p(p-2)/(p-1)^2 > 0$: Jeder Faktor ist positiv
      (für $p \geq 3$: $p-2 \geq 1 > 0$), das Produkt konvergiert absolut
      gegen $C_2 \approx 0{,}6602 > 0$.
    \item $(p-1)/(p-2) > 0$ für alle $p \geq 3$.
    \item Faktor $2 > 0$.
  \end{itemize}
  Also $\mathfrak{S}(n) > 0$ für alle geraden $n \geq 4$.
\end{proof}

\begin{bemerkung}[Ungerades $n$]
  Für ungerades $n$: $c_2(n) = -1$, Faktor bei $p=2$: $1+(-1)/1 = 0$.
  Daher $\mathfrak{S}(n) = 0$ für ungerades $n$.  Dies spiegelt korrekt wider,
  dass zwei ungerade Primzahlen eine gerade Summe ergeben.
\end{bemerkung}

% ================================================================
\section{Numerische Werte von \texorpdfstring{$\mathfrak{S}(n)$}{S(n)}}
\label{sec:numerik}
% ================================================================

\begin{center}
\begin{tabular}{rccrl}
\toprule
$n$ & $\mathfrak{S}(n)$ (exakt) & $\approx$ & $r_2(n)$ & Beispiele \\
\midrule
4  & $2C_2$ & $1{,}320$ & 1 & $2+2$ \\
6  & $2C_2 \cdot \frac{p-1}{p-2}\big|_{p=3} = 2C_2 \cdot \frac{2}{1} = 4C_2$ & $2{,}641$ & 1 & $3+3$ \\
8  & $2C_2$ & $1{,}320$ & 1 & $3+5$ \\
10 & $2C_2 \cdot \frac{4}{3} = \frac{8}{3}C_2$ & $1{,}760$ & 2 & $3+7,\, 5+5$ \\
12 & $2C_2 \cdot \frac{2}{1} = 4C_2$ & $2{,}641$ & 1 & $5+7$ \\
30 & $2C_2 \cdot \frac{2}{1}\cdot\frac{4}{3} = \frac{16}{3}C_2$
   & $3{,}521$ & 3 & $7+23,\, 11+19,\, 13+17$ \\
\bottomrule
\end{tabular}
\end{center}

% ================================================================
\section{Die Zwillingsprimzahl-Konstante und \texorpdfstring{$C_2$}{C2}}
\label{sec:zwillingsprim}
% ================================================================

Die Konstante $C_2$ in $\mathfrak{S}(n)$ ist dieselbe wie die
\emph{Zwillingsprimzahl-Konstante}:

\begin{theorem}[Hardy--Littlewood, 1923]\label{thm:zp}
  Bedingt auf Vermutung~$\mathbf{H}$ (Abschnitt~\ref{sec:ktuples}):
  \[
    \pi_2(x) \;\sim\; 2 C_2 \cdot \frac{x}{(\log x)^2},
  \]
  wobei $C_2 \approx 0{,}6601618\ldots$
\end{theorem}

Numerische Annäherung:
\[
  C_2 \;=\; \frac{15}{16} \cdot \frac{35}{36} \cdot \frac{143}{144} \cdots
  \;=\; \prod_{p \geq 3} \left(1 - \frac{1}{(p-1)^2}\right)
  \;\approx\; 0{,}66016182\ldots
\]

% ================================================================
\section{Die Hardy--Littlewood-Vermutung~$\mathbf{H}$\\
         über Primzahl-$k$-Tupel}
\label{sec:ktuples}
% ================================================================

\begin{definition}[Zulässiges $k$-Tupel]
  Ein $k$-Tupel $\mathcal{H} = (h_1, \ldots, h_k)$ nicht-negativer ganzer Zahlen
  heißt \emph{zulässig}, wenn für jede Primzahl $p$:
  \[
    \bigl|\{h_i \pmod p : 1 \leq i \leq k\}\bigr| < p.
  \]
\end{definition}

\begin{bemerkung}
  $\{0, 2\}$ (Zwillingsprimzahlen): zulässig ($|\{0,2\} \pmod p| \leq 2 < p$
  für alle $p$).  $\{0, 1, 2, 3, 4, 5\}$: nicht zulässig für $p = 3$.
\end{bemerkung}

\begin{vermutung}[Hardy--Littlewood Hypothese~$\mathbf{H}$, 1923]\label{verm:H}
  Für jedes zulässige $k$-Tupel $\mathcal{H} = (h_1, \ldots, h_k)$ gilt:
  \[
    \pi_{\mathcal{H}}(x)
    \;=\; \#\bigl\{n \leq x : n+h_1, \ldots, n+h_k \text{ alle prim}\bigr\}
    \;\sim\; \mathfrak{S}(\mathcal{H}) \cdot \frac{x}{(\log x)^k},
  \]
  wobei
  \[
    \mathfrak{S}(\mathcal{H})
    \;=\; \prod_p \left(1-\frac{1}{p}\right)^{-k} \left(1-\frac{\nu_p(\mathcal{H})}{p}\right)
  \]
  und $\nu_p(\mathcal{H}) = |\{h_i \pmod p\}|$.
\end{vermutung}

\begin{beispiel}[Spezialfälle von Vermutung~$\mathbf{H}$]
  \begin{enumerate}
    \item $k=1$, $\mathcal{H}=\{0\}$:
      $\pi(x) \sim x/\log x$ (Primzahlsatz ✓).
    \item $k=2$, $\mathcal{H}=\{0,2\}$:
      $\pi_2(x) \sim 2C_2 \cdot x/(\log x)^2$ (Zwillingsprimzahl-Vermutung).
    \item $k=2$, $\mathcal{H}=\{0,4\}$ (Cousin-Primzahlen $(p,p+4)$):
      dieselbe Konstante $2C_2$!
    \item $k=3$, $\mathcal{H}=\{0,2,6\}$ (Primzahldrillinge):
      $\mathfrak{S}(\mathcal{H}) = 9 \prod_{p \geq 5} p^2(p-3)/(p-1)^3$.
    \item Goldbach ($k=2$): $n = p+q$ entspricht $\pi_{\{0,n\}}(N)$ für $N=n$,
      mit $\mathfrak{S}(\{0,n\}) = \mathfrak{S}(n)$ wie oben.
  \end{enumerate}
\end{beispiel}

\begin{bemerkung}[Zulässigkeit und Goldbach-Vermutung]
  Das Tupel $\{0, n\}$ ist für gerades $n$ zulässig:
  $|\{0 \pmod p, n \pmod p\}| \leq 2 < p$ für $p \geq 3$;
  für $p=2$: $\{0 \pmod 2, n \pmod 2\} = \{0\}$ (da $n$ gerade), also $|\{0\}|=1<2$.
  Daher sagt Vermutung~$\mathbf{H}$ eine positive Darstellungsanzahl für
  alle geraden $n$ vorher — konsistent mit Goldbach.

  Für ungerades $n$: $\{0 \pmod 2, n \pmod 2\} = \{0,1\}$ deckt beide
  Restklassen mod $2$ ab, also ist $\mathfrak{S}(\{0,n\}) = 0$ und das
  Tupel ist \emph{nicht} zulässig — konsistent damit, dass zwei Primzahlen
  keine ungerade Zahl $\geq 5$ ergeben können.
\end{bemerkung}

% ================================================================
\section{Das Paritätsproblem und die Grenzen der Kreismethode}
\label{sec:paritaet}
% ================================================================

\begin{bemerkung}[Warum die Kreismethode binäres Goldbach nicht beweisen kann]
  Der Nebenbogen-Fehler für $\int_{\mathfrak{m}} S(\alpha)^2 \eq{-n\alpha}\, d\alpha$
  ist:
  \[
    \left|\int_{\mathfrak{m}} S^2 \eq{-n\cdot}\right|
    \;\leq\; \|S\|_{\mathfrak{m},\infty} \cdot \int_0^1 |S|^2
    \;\ll\; N(\log N)^{-B} \cdot N\log N
    \;=\; N^2 (\log N)^{1-B}.
  \]
  Der Hauptterm aus Hauptbögen ist $\Theta(N)$ (für $n = N$).  Da
  $N^2 (\log N)^{1-B} \gg N$ für alle Wahlen von $B$, überwiegt
  der Fehler den Hauptterm.

  Fundamentale Erklärung: Das \emph{Paritätsproblem} (Selberg, 1949):
  Siebmethoden und die Kreismethode können Primzahlen ($\mathbb{P}_1$)
  nicht von Produkten zweier Primzahlen ($\mathbb{P}_2$) unterscheiden,
  da beide „ungerade" sind in einem Paritätssinn.
\end{bemerkung}

\begin{theorem}[Formel für $R_2(n)$ unter GRH]\label{thm:R2_GRH}
  Unter der Verallgemeinerten Riemannschen Hypothese (GRH):
  \[
    R_2(n) \;=\; \mathfrak{S}(n) \cdot n \;+\; O\!\left(\sqrt{n}\, (\log n)^2\right).
  \]
  Insbesondere gilt $R_2(n) > 0$ für alle hinreichend großen geraden $n$
  (binäres Goldbach für große $n$ unter GRH).
\end{theorem}

% ================================================================
\section{Numerik und der Goldbach-Komet}
\label{sec:komet}
% ================================================================

Die Funktion $r_2(n)$ für gerades $n$, aufgetragen gegen $n$, bildet den
berühmten \emph{Goldbach-Kometen}: ein Streudiagramm von $(n, r_2(n))$, das
der Kurve $\mathfrak{S}(n) \cdot n / (\log n)^2$ eng folgt.

\begin{bemerkung}[Struktur des Goldbach-Kometen]
  Der Graph von $r_2(n)$ (Anzahl der Goldbach-Zerlegungen von $n$) gegen
  $n$ heißt der \emph{Goldbach-Komet}.  Die Struktur:
  \begin{itemize}
    \item Gerades $n$ mit vielen kleinen Primteilern hat größeres
      $\mathfrak{S}(n)$ und damit mehr Darstellungen.
    \item Primoriale $n = 2 \cdot 3 \cdot 5 \cdots p_k$ geben das
      maximale $\mathfrak{S}(n)$ unter Zahlen vergleichbarer Größe.
    \item $n = 2^k$ hat minimales $\mathfrak{S}(n) = 2C_2$ (keine
      ungeraden Primteiler).
  \end{itemize}
\end{bemerkung}

\begin{center}
\begin{tabular}{rlll}
\toprule
$n$ & $\mathfrak{S}(n)$ & Prognose $r_2(n)$ & Tatsächlich $r_2(n)$ \\
\midrule
100    & $2C_2 \cdot \tfrac{4}{3} \approx 1{,}76$ & $\approx 6$ & 6 \\
1000   & variabel & $\approx 18$ & 28 \\
$10^6$ & variabel & $\approx 1285$ & 1302 \\
$10^{10}$ & variabel & $\approx 1{,}5\times 10^7$ & $\approx 1{,}5\times 10^7$ \\
\bottomrule
\end{tabular}
\end{center}

% ================================================================
\section{Bezug zu Chens Satz}
\label{sec:chen}
% ================================================================

\begin{theorem}[Chen Jingrun, 1973 — vgl.\ Paper~16]\label{thm:chen}
  Jede hinreichend große gerade Zahl $n$ ist darstellbar als $n = p + m$,
  wobei $p$ Primzahl und $m$ Produkt von höchstens zwei Primzahlen ist
  ($m \in \mathbb{P}_2$).
\end{theorem}

\begin{bemerkung}[Paritätslücke: Chen vs.\ Goldbach]
  Chens Satz sagt $n = p + P_2$, nicht $n = p + p$.  Der Grund:
  Das Paritätsproblem verhindert die Verfeinerung von $P_2$ auf $P_1$
  (echte Primzahlen).  Die singuläre Reihe $\mathfrak{S}(n)$ erscheint
  in beiden Vorhersagen mit demselben Ausdruck, aber unterschiedlichem
  Vorfaktor:
  \[
    \#\{p : p + P_2 = n\} \;\sim\; C \cdot \mathfrak{S}(n) \cdot \frac{n}{(\log n)^2}.
  \]
\end{bemerkung}

% ================================================================
\section{Schlussfolgerung}
% ================================================================

Die Goldbach'sche singuläre Reihe $\mathfrak{S}(n)$ und die Hardy--Littlewood-Vermutung~$\mathbf{H}$
repräsentieren unser bestes quantitatives Verständnis des binären Goldbach-Problems.

Hauptergebnisse dieses Papers:
\begin{enumerate}
  \item $\mathfrak{S}(n) = 2C_2 \prod_{p|n, p\geq 3}(p-1)/(p-2) > 0$
    für alle geraden $n \geq 4$ (Sätze~\ref{thm:euler},~\ref{thm:pos}).
  \item Unter GRH: binäres Goldbach gilt für alle hinreichend großen
    geraden $n$ (Satz~\ref{thm:R2_GRH}).
  \item Vermutung~$\mathbf{H}$ vereinheitlicht alle Primzahlkonstellation-Probleme.
  \item Das Paritätsproblem erklärt, warum die Kreismethode allein das
    binäre Goldbach-Problem nicht unbedingt beweisen kann.
\end{enumerate}

Die binäre Goldbach-Vermutung und die Zwillingsprimzahl-Vermutung gehören
zu den tiefsten offenen Problemen der Mathematik.  Fortschritte erfordern
das Überwinden der Paritätsbarriere in der Siebtheorie.

% ================================================================
\begin{thebibliography}{9}

\bibitem{HardyLittlewood1923}
G.~H.~Hardy und J.~E.~Littlewood.
\newblock Some problems of ``Partitio Numerorum''.  III.
\newblock \emph{Acta Math.}, 44:1--70, 1923.

\bibitem{Helfgott2013a}
H.~A.~Helfgott.
\newblock Major arcs for Goldbach's theorem.
\newblock arXiv:1305.2897, 2013.

\bibitem{Chen1973}
C.~J.~Chen.
\newblock On the representation of a large even integer as the sum of a prime
  and the product of at most two primes.
\newblock \emph{Sci.\ Sinica}, 16:157--176, 1973.

\bibitem{Vaughan1981}
R.~C.~Vaughan.
\newblock \emph{The Hardy--Littlewood Method}, 2.~Aufl.
\newblock Cambridge University Press, 1997.

\bibitem{IwaniecKowalski}
H.~Iwaniec und E.~Kowalski.
\newblock \emph{Analytic Number Theory}.
\newblock American Mathematical Society, 2004.

\bibitem{OliveiraSilva2014}
T.~Oliveira~e~Silva, S.~Herzog und S.~Pardi.
\newblock Empirical verification of the even Goldbach conjecture.
\newblock \emph{Math.\ Comp.}, 83(288):2033--2060, 2014.

\bibitem{Goldston1992}
D.~A.~Goldston.
\newblock On Hardy and Littlewood's contribution to the Goldbach conjecture.
\newblock In \emph{Proceedings of the Amalfi Conference on Analytic Number
  Theory}, pp.\ 115--155.  Universit\`a di Salerno, 1992.

\bibitem{Vinogradov1937}
I.~M.~Vinogradov.
\newblock Representation of an odd number as a sum of three primes.
\newblock \emph{Dokl.\ Akad.\ Nauk SSSR}, 15:169--172, 1937.

\bibitem{Granville1995}
A.~Granville.
\newblock Harald Cram\'er and the distribution of prime numbers.
\newblock \emph{Scand.\ Actuarial J.}, 1:12--28, 1995.

\end{thebibliography}

\end{document}
